\documentclass[10pt]{article}
\usepackage[T1]{fontenc}
\usepackage[utf8]{inputenc}
\usepackage[french]{babel}
\usepackage{amsmath,amssymb,mathrsfs,mathpazo}
\usepackage{geometry}
\geometry{a4paper,margin=2cm}
\usepackage{graphicx,array,tabularx,tabularray,colortbl,multirow,multicol,booktabs,enumitem,pifont,ulem,fancybox,diagbox,cancel}
\usepackage{pstricks,pstricks-add,pst-plot,pst-tree,pst-node,pst-text,pst-func,pst-math,pst-bspline,pst-3dplot,pst-eucl,pst-solides3d}
\let\pscalculate\undefined
\usepackage{pst-bezier}
\usepackage{tikz,pgfplots,tkz-tab}
\usetikzlibrary{arrows,arrows.meta,patterns,calc,positioning,trees,decorations.markings,shapes.geometric,shapes.misc,shapes.arrows,intersections,decorations.pathmorphing,decorations.pathreplacing,backgrounds,shadows,babel}
\pgfplotsset{compat=1.18}
\usepgfplotslibrary{fillbetween}
\usepackage[np]{numprint}
\usepackage{listings,makecell,pas-tableur,twoopt}
\let\barre\undefined
\usepackage[pstricks]{tablvar}
\definecolor{darkgreen}{RGB}{0,100,0}
\providecommand{\red}{\color{red}}
\providecommand{\blue}{\color{blue}}
\providecommand{\green}{\color{green}}
\providecommand{\black}{\color{black}}
\providecommand{\white}{\color{white}}
\providecommand{\gray}{\color{gray}}
\providecommand{\cyan}{\color{cyan}}
\providecommand{\magenta}{\color{magenta}}
\providecommand{\yellow}{\color{yellow}}
\providecommand{\violet}{\color{violet}}
\providecommand{\orange}{\color{orange}}
\providecommand{\euro}{\texteuro}
\providecommand{\EUR}{\texteuro}
\providecommand{\ieme}{\textsuperscript{e}}
\providecommand{\degres}{\ensuremath{^\circ}}
\providecommand{\degree}{\ensuremath{^\circ}}
\providecommand{\vv}[1]{\overrightarrow{#1}}
\providecommand{\Coord}[1]{#1}
\providecommand{\decosix}{}
\providecommand{\starredbullet}{\ensuremath{\bullet}}
\pagestyle{empty}
\setlength{\parindent}{0pt}
\usepackage{adjustbox}\DeclareUnicodeCharacter{2236}{:}\DeclareUnicodeCharacter{2190}{\ensuremath{\leftarrow}}
\providecommand{\R}{\mathbb{R}}
\providecommand{\N}{\mathbb{N}}
\providecommand{\Z}{\mathbb{Z}}
\providecommand{\Q}{\mathbb{Q}}
\providecommand{\D}{\mathbb{D}}
\providecommand{\C}{\mathbb{C}}
\providecommand{\e}{{\mathrm{e}}}
\providecommand{\vect}[1]{\overrightarrow{#1}}
\providecommand{\vectt}[1]{\overrightarrow{\mathrm{#1}}}
\providecommand{\vvt}[1]{\overrightarrow{\mathrm{#1}}}
\providecommand{\np}[2][]{{#2~#1}}
\providecommand{\vv}[1]{{\overrightarrow{#1}}}
\renewcommand{\d}{{\mathrm{d}}}
\providecommand{\degres}{{^\circ}}
\providecommand{\degre}{{^\circ}}
\providecommand{\blue}{\color{blue}}
\providecommand{\red}{\color{red}}
\providecommand{\enskip}{\quad}
\providecommand{\llbracket}{\lbrack\!\lbrack}
\providecommand{\rrbracket}{\rbrack\!\rbrack}
\providecommand{\euro}{{\text{€}}}
\providecommand{\no}{{\text{n°}}}
\providecommand{\interval}[2]{{[#1;#2]}}
\providecommand{\intervFO}[2]{{[#1;#2[}}
\providecommand{\dotfill}{\ldots}

\providecommand{\ds}{\displaystyle}
\providecommand{\textding}[1]{#1}
\providecommand{\Oij}{\(\mathrm{O};\vec{\imath},\vec{\jmath}\)}
\providecommand{\Oijk}{\(\mathrm{O};\vec{\imath},\vec{\jmath},\vec{k}\)}
\providecommand{\Ouv}{\(\mathrm{O};\vec{u},\vec{v}\)}
\providecommand{\intervFF}[2]{\ensuremath{[#1;#2]}}
\renewcommand{\euro}{\ifmmode\text{\texteuro}\else\texteuro\fi}
\DeclareUnicodeCharacter{00B0}{\ifmmode^{\circ}\else\textdegree\fi}
\DeclareUnicodeCharacter{20AC}{\ifmmode\text{\texteuro}\else\texteuro\fi}
% Énoncés : Éducation nationale. Transcription et corrigés : APMEP, Denis Vergès et contributeurs. Découpage et classement : CoopMaths, Sébastien Lozano, Éric Elter et contributeurs. Adaptation de la présentation : Maths Explorer.
% CC BY-SA : https://coopmaths.fr/www/about/
% https://coopmaths.fr/alea/static/bac/2022/tex/bac_2022_10_sujet2_caledonie_1.tex
\begin{document}
\noindent\textit{APMEP / CoopMaths. CC BY-SA. Présentation adaptée par Maths Explorer.}\par\medskip
\textbf{Principaux domaines abordés :} Probabilités

\medskip

Au basket-ball, il existe deux sortes de tir :

\begin{itemize}
\item les tirs à deux points.

Ils sont réalisés près du panier et rapportent deux points s'ils sont réussis.
\item les tirs à trois points.

Ils sont réalisés loin du panier et rapportent trois points s'ils sont réussis.
\end{itemize}

\smallskip

Stéphanie s'entraîne au tir. On dispose des données suivantes :

\begin{itemize}
\item[$\bullet~~$] Un quart de ses tirs sont des tirs à deux points. Parmi eux, 60\,\% sont réussis.
\item[$\bullet~~$] Trois quarts de ses tirs sont des tirs à trois points. Parmi eux, 35\,\% sont réussis.
\end{itemize}

\medskip

\begin{enumerate}
\item Stéphanie réalise un tir. 

On considère les évènements suivants :

\begin{description}
\item[ ] $D$ : \og Il s'agit d'un tir à deux points \fg.
\item[ ] $R$ : \og le tir est réussi \fg.
\end{description}
	\begin{enumerate}
		\item Représenter la situation à l'aide d'un arbre de probabilités.
		\item Calculer la probabilité $p\left(\overline{D} \cap R\right)$.
		\item Démontrer que la probabilité que Stéphanie réussisse un tir est égale à $0,4125$.
		\item Stéphanie réussit un tir. Calculer la probabilité qu'il s'agisse d'un tir à trois points. Arrondir le résultat au centième.
	\end{enumerate}	
\item Stéphanie réalise à présent une série de $10$ tirs à trois points.

On note $X$ la variable aléatoire qui compte le nombre de tirs réussis.

On considère que les tirs sont indépendants. On rappelle que la probabilité que Stéphanie réussisse un tir à trois points est égale à $0,35$.
	\begin{enumerate}
		\item Justifier que $X$ suit une loi binomiale. Préciser ses paramètres.
		\item Calculer l'espérance de $X$. Interpréter le résultat dans le contexte de l'exercice.
		\item Déterminer la probabilité que Stéphanie rate $4$ tirs ou plus. Arrondir le résultat au centième.
		\item Déterminer la probabilité que Stéphanie rate au plus $4$ tirs. Arrondir le résultat au centième.
	\end{enumerate}
\item Soit $n$ un entier naturel non nul.

Stéphanie souhaite réaliser une série de $n$ tirs à trois points. 

On considère que les tirs sont indépendants. On rappelle que la probabilité qu'elle réussisse un tir à trois points est égale à $0,35$.

Déterminer la valeur minimale de $n$ pour que la probabilité que Stéphanie réussisse au moins un tir parmi les $n$ tirs soit supérieure ou égale à $0,99$.
\end{enumerate}

\bigskip

\end{document}
